%! AP Statistics — Unit 7, Lesson 7.1 — Student Packet Cover Sheet
\documentclass[10pt]{article}
\usepackage{apstats-article}
\usepackage{apstats-boxes}
\usepackage{tikz}

\begin{document}

\begin{tikzpicture}[remember picture, overlay]
  \fill[navy] (current page.north west) rectangle
    ([yshift=-2.2cm]current page.north east);
  \node[anchor=west, white, font=\large\bfseries] at
    ([xshift=1cm, yshift=-0.7cm]current page.north west)
    {\CourseName: \SchoolYear};
  \node[anchor=west, white, font=\normalsize] at
    ([xshift=1cm, yshift=-1.3cm]current page.north west)
    {Unit 7: Inference for Quantitative Data: Means};
  \node[anchor=west, white, font=\normalsize\bfseries] at
    ([xshift=1cm, yshift=-1.9cm]current page.north west)
    {Lesson 7.1 \quad Introducing Statistics: Why Should I Worry About Error?};
\end{tikzpicture}

\vspace{2.4cm}
\namedateperiod

\begin{learningtargetbox}
\begin{itemize}
    \item I can describe what it means to commit a Type I error in the context of a
          statistical inference problem.
    \item I can describe what it means to commit a Type II error in the context of a
          statistical inference problem.
    \item I can explain how the significance level $\alpha$ controls the probability
          of a Type I error and how changing $\alpha$ affects the Type II error rate.
\end{itemize}
\end{learningtargetbox}

\begin{tocbox}
\renewcommand{\arraystretch}{1.5}
\begin{tabularx}{\linewidth}{clXc}
    \textbf{\#} & \textbf{Component} & \textbf{Description} & \textbf{Score} \\
    \hline
    1 & Warm-Up      & Spiral review: sampling distributions and error types & \hfill/5 \\
    2 & Guided Notes & Vocabulary, error types, framing Unit 7               & \hfill/10 \\
    3 & Activity     & Partner scenarios: identifying and evaluating errors   & \hfill/12 \\
    4 & Exit Ticket  & Independent check on Type I/II errors and $\alpha$     & \hfill/6 \\
    5 & Homework     & Practice with error types across multiple contexts     & \hfill/12 \\
    \hline
       & \multicolumn{2}{l}{\textbf{Total}} & \hfill/45 \\
\end{tabularx}
\end{tocbox}

\end{document}
